% ---------------------------------------------------------------------------
%  APPENDICES
% ---------------------------------------------------------------------------
\chapter{Past-Paper Index --- Where Each Question Has Appeared}
\chaptermark{Past-Paper Index}

The five papers with surviving scans are reproduced in full in Part~III. This
appendix indexes the \textbf{other four} question by question, with the same
pointers to the answers --- \textbf{Q\emph{n}} for Part~I, \textbf{A\emph{n}}
\,/\, \textbf{B\emph{n}} \,/\, \textbf{C\emph{n}} for the three sections of
Part~II. It closes with a one-line summary of all nine papers.

\section*{A.1\enspace CIE 1 --- 22 June 2024 \ \ \normalfont\small(HS248AT,
  25 marks)}
\addcontentsline{toc}{section}{A.1 CIE 1, 22 June 2024}

\begin{idxtable}
\qn{1} & Define self-exploration. What is the content of self-exploration? &
  5 & \ap{Q8} \\ \hline
\qn{2} & ``Right understanding $+$ Relationship $=$ Mutual fulfilment; Right
  understanding $+$ Physical facilities $=$ Mutual prosperity.'' Illustrate
  with two examples for each. & 8 & \ap{Q20} \\ \hline
\qn{3} & How can we ensure harmony in self (`I')? & 4 & \ap{Q24} \\ \hline
\qn{4} & How do prosperity and wealth differ? What is more acceptable to us
  and why? & 8 & \ap{Q16} \\ \hline
\end{idxtable}

\section*{A.2\enspace CIE 2 --- 24 July 2024 \ \ \normalfont\small(HS248AT,
  5 $+$ 25 marks)}
\addcontentsline{toc}{section}{A.2 CIE 2, 24 July 2024}

\begin{idxtable}
\qn{A1--5} & Quiz: foundational values; undivided society; comprehensive human
  goals; gratitude & 5 & \ap{A6--A16} \\ \hline
\qn{B1} & What is the difference between respect and disrespect? Which of the
  two is naturally acceptable to you? & 5 & \ap{Q37} \\ \hline
\qn{B2} & How does affection lead to harmony in the family? & 5 & \ap{Q38}
  \\ \hline
\qn{B3} & What is the meaning of justice in human relationships? How does it
  follow from family to world family? & 5 & \ap{Q29} \\ \hline
\qn{B4} & What do you understand by trust? Differentiate between intention and
  competence with examples. & 5 & \ap{Q32} \\ \hline
\qn{B5} & Explain the comprehensive human goal. How does fearlessness follow
  from right understanding and prosperity? & 5 & \ap{Q42} \\ \hline
\end{idxtable}

\section*{A.3\enspace SEE --- September/October 2024 and October 2023
  \ \ \normalfont\small(50 marks)}
\addcontentsline{toc}{section}{A.3 SEE, Sept/Oct 2024 and Oct 2023}

Where a slot carries two questions, the second is the alternative set in the
other year's paper. Both are answered.

\begin{idxtable}
\qn{1.1--1.10} & Part A --- objective & 10 & \ap{A1--A16} \\ \hline
\qn{2a} & ``To be in a state of harmony is happiness.'' Explain and illustrate
  with two examples. & 6 & \ap{Q14} \\ \hline
\qn{2a\,alt} & Explain the process of self-exploration with a diagram. & 6 &
  \ap{Q9} \\ \hline
\qn{2b} & What do you mean by values? How do they differ from skills? How are
  values and skills complementary? & 6 & \ap{Q6} \\ \hline
\qn{2b\,alt} & What is value education? Describe the need for value education
  with its guidelines. & 6 & \ap{Q1}\newline\ap{Q2}\newline\ap{Q3} \\ \hline
\qn{3a} & What is the role of the value system in family harmony? & 5 &
  \ap{Q48} \\ \hline
\qn{3a\,alt} & Explain the problems faced due to differentiation in
  relationship. & 5 & \ap{Q36} \\ \hline
\qn{3b} & ``Discrimination leads to acrimony in relationships.'' Explain. What
  problems are created when we discriminate? & 5 & \ap{Q36} \\ \hline
\qn{3b\,alt} & What is `justice'? What are its four elements? Is it a
  continuous or a temporary need? & 5 & \ap{Q28} \\ \hline
\qn{3c} & Define love. How can you say that love is the complete value? & 4 &
  \ap{Q40} \\ \hline
\qn{4a} & What do you understand by trust? Differentiate between intention and
  competence with examples. & 5 & \ap{Q32} \\ \hline
\qn{4a\,alt} & Distinguish between the needs of the Self and the needs of the
  Body. & 5 & \ap{Q23} \\ \hline
\qn{4b} & Realization and understanding are essential for happiness and
  harmony. Explain. & 5 & \ap{Q26} \\ \hline
\qn{4b\,alt} & The feeling of love lays down the basis of an undivided
  society. Explain. & 5 & \ap{Q40} \\ \hline
\qn{4c} & Describe the concept of an undivided society and the universal
  order, and explain how both can help create a world family. & 5 & \ap{Q46}
  \\ \hline
\qn{4c\,alt} & Bring out the differences between respect and differentiation.
  & 4 & \ap{Q37} \\ \hline
\qn{5a} & Existence is co-existence of mutually interacting units in
  all-pervasive space. Explain. \,/\, Describe briefly the nature of
  co-existence. & 5 & \ap{Q62}\newline\ap{Q63} \\ \hline
\qn{5b} & Define harmony in nature and why is it important. Explain with
  examples. & 5 & \ap{Q50} \\ \hline
\qn{5b\,alt} & What are the four orders of nature? Briefly explain them. & 5 &
  \ap{Q51} \\ \hline
\qn{5c} & What do you mean by co-existence? & 4 & \ap{Q63} \\ \hline
\qn{6a} & What is the natural characteristic (swabhava) of the human order?
  Explain. & 5 & \ap{Q54} \\ \hline
\qn{6b} & What are the four orders of nature? Briefly explain them. & 5 &
  \ap{Q51} \\ \hline
\qn{6c} & How can we say that `nature is Self-Organized and in space
  Self-Organization is Available'? & 4 & \ap{Q61} \\ \hline
\qn{5c\,/\,6b}\newline\qn{(2023)} & \emph{Holistic technology}; \emph{ethical
  human conduct in terms of values, policy and character}; \emph{nature of the
  universal human order} & 4--5 & --- \\ \hline
\end{idxtable}

\begin{note}
\textbf{\sffamily The one gap in this book, stated plainly.} The last row
above belongs to \textbf{Chapters 11--13}, which are \emph{not} part of the
Unit I--III material supplied for this course. \emph{Holistic technology} and
\emph{ethical human conduct in terms of values, policy and character} are
therefore \textbf{not answered here} --- no source was provided for them, and
answering from outside the prescribed material would break the rule this book
is written under. (The third of the three, the \emph{nature of the universal
human order}, is within scope and is answered in \textbf{Q46}.) If those two
topics are in your syllabus this year, get them from the Chapter 11--13
notes.
\end{note}

\section*{A.4\enspace All nine papers at a glance}
\addcontentsline{toc}{section}{A.4 All nine papers at a glance}

\renewcommand{\arraystretch}{1.35}
\arrayrulecolor{rulegrey}
\begin{tabular}{@{}P{44mm} P{30mm} C{18mm} C{22mm} P{34mm}@{}}
\rowcolor{slate}
\hdr{Assessment} & \hdr{Date} & \hdr{Max M} & \hdr{Questions} & \hdr{Where}
  \\ \hline
SEE (21HSU48) & October 2023 & 50 & 23 & Appendix A.3 \\ \hline
CIE 1 & 22 June 2024 & 25 & 4 & Appendix A.1 \\ \hline
CIE 2 & 24 July 2024 & 5 $+$ 25 & 10 & Appendix A.2 \\ \hline
Improvement CIE \,/\, CIE 3 & 28 August 2024 & 25 & 10 &
  \textbf{Part III, \S1} \\ \hline
SEE & Sept/Oct 2024 & 50 & 23 & Appendix A.3 \\ \hline
CIE-II & 7 May 2025 & 25 $+$ 5 & 10 & \textbf{Part III, \S2} \\ \hline
Improvement CIE & June 2025 & 30 & 9 & \textbf{Part III, \S3} \\ \hline
\textbf{SEE} Regular \,/\, Suppl. & June/July 2025 & 50 & 20 &
  \textbf{Part III, \S4} \\ \hline
Improvement CIE & 18 June 2026 & 05 $+$ 25 & 10 & \textbf{Part III, \S5}
  \\ \hline
\end{tabular}

\chapter{Questions Set More Than Once}
\chaptermark{Questions Set More Than Once}

Across the nine papers, ten topics have been examined two, three or four
times. They are listed below with every occurrence. \textbf{This is the single
most useful page in the book if revision time is short} --- between them these
ten topics account for well over a third of every mark ever set on this
course.

\vspace{2pt}
Watch what happens to the marks and the Bloom's level as a question moves
between papers. \emph{Conformance in the four orders} was worth 6 marks on a
CIE and 7 on the SEE, at L4 both times: the semester-end paper wanted the same
analysis, one point longer. That is the general pattern --- a repeat keeps its
level and gains a mark or two on the way to the SEE. The exception is the
\textbf{18 June 2026} paper, which took three long-answer staples and set them
as \emph{one-mark} short answers; the same material, compressed to a single
line.

\renewcommand{\arraystretch}{1.35}
\arrayrulecolor{rulegrey}
\setlength{\LTpre}{8pt}\setlength{\LTpost}{4pt}
\begin{longtable}{P{58mm} P{80mm} C{14mm}}
\rowcolor{slate}
\hdr{Topic} & \hdr{Set in} & \hdr{Answer} \\ \hline \endfirsthead
\rowcolor{slate}
\hdr{Topic} & \hdr{Set in} & \hdr{Answer} \\ \hline \endhead
\textbf{The four orders of nature} --- set five times, at four different
  lengths: a five-mark explanation, a six-mark version with a schematic
  diagram, and a one-mark list &
  Improvement CIE, 28 Aug 2024, Part B Q3 (5 M)\newline
  SEE, Sept/Oct 2024, Q5b\,alt and Q6b (5 M)\newline
  Improvement CIE, Jun 2025, Test Q1 (6 M)\newline
  Improvement CIE, 18 Jun 2026, Part A Q2 (1 M) &
  \ap{Q51}\newline\ap{Q74}\newline\ap{C2} \\ \hline
\textbf{Svabhava --- the natural characteristic} --- narrowed and widened
  across four papers: the human order alone, then all four orders, then the
  self &
  Improvement CIE, 28 Aug 2024, Part B Q2 (5 M)\newline
  SEE, Sept/Oct 2024, Q6a (5 M)\newline
  SEE, Jun/Jul 2025, Q5b (07 M)\newline
  Improvement CIE, 18 Jun 2026, Part A Q3 (1 M) &
  \ap{Q54}\newline\ap{C3} \\ \hline
\textbf{Co-existence} --- the four-mark definition, the five-mark
  ``mutually interacting units'' version, and the one-mark ``two
  definitions'' &
  Improvement CIE, 28 Aug 2024, Part B Q1 (5 M)\newline
  SEE, Sept/Oct 2024, Q5a and Q5c (5 M, 4 M)\newline
  SEE, Jun/Jul 2025, Q6b (06 M)\newline
  Improvement CIE, 18 Jun 2026, Part A Q1 (1 M) &
  \ap{Q62}\newline\ap{Q63}\newline\ap{C1} \\ \hline
\textbf{`Conformance', and conformance in the four orders} --- identical
  wording twice, plus the \emph{sanskaar} variant &
  Improvement CIE, Jun 2025, Test Q2 (6 M, CO2, L4)\newline
  SEE, Jun/Jul 2025, Q5a --- ``What is sanskaar?'' (06 M, L4)\newline
  SEE, Jun/Jul 2025, Q6a (07 M, CO3, L4) &
  \ap{Q55} \\ \hline
\textbf{Trust --- intention versus competence} --- identical wording all
  three times &
  CIE 2, 24 Jul 2024, Q B4 (5 M)\newline
  SEE, Sept/Oct 2024, Q4a (5 M)\newline
  CIE-II, 7 May 2025, Part B Q2 (5 M) &
  \ap{Q32} \\ \hline
\textbf{Differentiation and discrimination in relationship} &
  CIE 2 \,/\, SEE, Q3a\,alt and Q3b (5 M)\newline
  CIE-II, 7 May 2025, Part B Q4 (5 M) &
  \ap{Q36} \\ \hline
\textbf{Units and space} --- a seven-mark elaboration, then a one-mark
  differentiation &
  Improvement CIE, Jun 2025, Test Q4 (7 M, CO4, L4)\newline
  Improvement CIE, 18 Jun 2026, Part A Q5 (1 M, CO4, L2) &
  \ap{Q65}\newline\ap{C5} \\ \hline
\textbf{Nature is self-organised; in space self-organisation is available} &
  Improvement CIE, 28 Aug 2024, Part B Q4 (5 M)\newline
  SEE, Sept/Oct 2024, Q6c (4 M) &
  \ap{Q61} \\ \hline
\textbf{Trust as the foundation of relationship} (objective) --- identical
  wording and identical options &
  CIE-II, 7 May 2025, Part A Q1 (1 M, CO1, L1)\newline
  SEE, Jun/Jul 2025, Q1.6 (01 M, CO1, L2) &
  \ap{A17}\newline\ap{A32} \\ \hline
\textbf{Utility value of physical facility} (objective) --- reworded:
  \emph{`in nurture, protection and providing means for the body'} became
  \emph{`to help and preserve its utility'}; the four options are unchanged &
  Improvement CIE, Jun 2025, Quiz Q1 (1 M, CO1, L2)\newline
  SEE, Jun/Jul 2025, Q1.8 (01 M, CO3, L2) &
  \ap{A22}\newline\ap{A34} \\ \hline
\textbf{Love as the complete value, and the undivided society} &
  SEE, Sept/Oct 2024, Q3c (4 M)\newline
  SEE, Sept/Oct 2024, Q4b\,alt (5 M) &
  \ap{Q40} \\ \hline
\textbf{Respect versus disrespect \,/\, differentiation} &
  CIE 2, 24 Jul 2024, Q B1 (5 M)\newline
  SEE, Sept/Oct 2024, Q4c\,alt (4 M) &
  \ap{Q37} \\ \hline
\end{longtable}

\chapter{Last-Minute Revision Priorities}
\chaptermark{Revision Priorities}

If you have one evening, work through this list and nothing else. Everything
on it has been set at least twice, or is the skeleton that several other
answers hang from.

\section*{C.1\enspace Unit I}
\addcontentsline{toc}{section}{C.1 Unit I}

\begin{abul}
\item The \textbf{five needs} for value education \seeq{2} and the
      \textbf{five guidelines} \seeq{3}.
\item The \textbf{process of self-exploration} --- the diagram, both
      mechanisms, and the three tests (assurance, satisfaction, universality)
      \seeq{9}.
\item The \textbf{five characteristics of natural acceptance} \seeq{13}.
\item \textbf{Prosperity versus wealth} --- the full comparison table
      \seeq{16}. Set for 8 marks twice.
\item \textbf{Right understanding $+$ relationship} and \textbf{right
      understanding $+$ physical facilities}, with two examples each
      \seeq{20}.
\item The \textbf{Self versus Body needs} table --- six rows \seeq{23}.
\item The \textbf{four-step process} for harmony in `I' \seeq{24}.
\end{abul}

\section*{C.2\enspace Unit II}
\addcontentsline{toc}{section}{C.2 Unit II}

\begin{abul}
\item \textbf{Justice and its four elements}, and that it is a
      \emph{continuous} need \seeq{28}. Follow the official scheme, not the
      legal categories --- see the warning in \emph{How This Book Is
      Arranged} if you have not read it.
\item \textbf{Intention versus competence}, with the friend-on-campus example
      \seeq{32}. Set three times.
\item \textbf{Respect as right evaluation}, and the \textbf{three mistakes}
      --- over, under, otherwise \seeq{35}.
\item The \textbf{differentiation table} --- body \,/\, physical facilities
      \,/\, beliefs, with the problems each creates \seeq{36}.
\item The \textbf{nine values} with their one-line definitions, and which one
      is the foundation value and which the complete value \seeq{31}.
\item The \textbf{comprehensive human goal} and its sequence: right
      understanding $\rightarrow$ prosperity $\rightarrow$ fearlessness
      $\rightarrow$ co-existence \seeq{41}.
\item The \textbf{five dimensions of human endeavour} and what each leads to
      \seeq{44}.
\end{abul}

\section*{C.3\enspace Unit III}
\addcontentsline{toc}{section}{C.3 Unit III}

\begin{abul}
\item \textbf{The single four-orders table} --- things, activity, innateness,
      svabhava, basic activity, conformance \seeq{52}. Most of Unit~III is a
      row or a column of this one table; if you memorise nothing else,
      memorise this.
\item \textbf{Existence $=$ space $+$ units}, and the five-row units-versus-
      space comparison \seeq{62}\seeq{65}.
\item \textbf{Self-organisation} --- units are self-organised; in space,
      self-organisation is \emph{available} \seeq{61}.
\item \textbf{Recyclability and self-regulation}, with the forest and breed
      examples \seeq{60}.
\item \textbf{Interconnectedness and mutual fulfilment} --- and the fact that
      the \emph{human order is the one missing link} \seeq{59}\seeq{74}.
\item \textbf{The two types of units} --- material (jada, gathansheel) and
      conscious (chaitanya, gathanpurna) \seeq{71}, and why development can
      only happen in the second \seeq{73}.
\end{abul}

\section*{C.4\enspace If you have ten minutes}
\addcontentsline{toc}{section}{C.4 If you have ten minutes}

\begin{keybox}
\textbf{Four definitions, verbatim.}
\begin{enumerate}[leftmargin=1.8em,itemsep=2pt,topsep=3pt]
\item \textbf{Trust} --- to be assured that each human being inherently wants
      oneself and the other to be happy and prosperous. \emph{The foundation
      value.}
\item \textbf{Respect} --- right evaluation; to be evaluated as I am.
\item \textbf{Justice} --- recognition of values, fulfilment, evaluation, and
      mutual happiness ensured.
\item \textbf{Prosperity} --- the feeling of having or making available
      \emph{more than the required} physical facilities.
\end{enumerate}

\vspace{3pt}
\textbf{Three sequences.}
\begin{enumerate}[leftmargin=1.8em,itemsep=2pt,topsep=3pt]
\item Right Understanding $\rightarrow$ Prosperity $\rightarrow$ Fearlessness
      $\rightarrow$ Co-existence \quad \emph{(comprehensive human goal)}
\item Realization \& Understanding $\rightarrow$ Desire $\rightarrow$ Thought
      $\rightarrow$ Expectation \quad \emph{(harmony in `I')}
\item Constitution $\rightarrow$ Seed $\rightarrow$ Breed $\rightarrow$
      Sanskar \quad \emph{(conformance in the four orders)}
\end{enumerate}

\vspace{3pt}
\textbf{One equation.} \quad Existence $=$ Space $+$ Units $=$ Nature
submerged in space.
\end{keybox}
